\chapter*{Preface}
\label{v4-ch:preface}

An equality in this volume is always an equality between two
constructions.  If the Borcherds--Kac--Moody central charge of a
frame-shape orbifold lift is written
\[
 \kappa_{\mathrm{BKM}}=\frac{c_N(0)}{2},
\]
then the formula is not tested by reading \(c_N(0)\) twice from the
same table.  It is tested by placing the Borcherds weight theorem on
one side and the Gritsenko--Nikulin denominator identity on the other.
A table may suggest the number; it cannot be the second proof of the
number.

This is the principle used throughout the volume.  To a theorem from
the first three volumes one attaches a comparison
\[
 \mathrm{Real}(X,T)=(d,e)
\]
whose two sides have disjoint mathematical origins and the same
invariant as their common value.  The symbol \(d\) records the
structural construction and \(e\) records the calculation.  The
intersection of the two arguments is the object compared, not the
method by which it is obtained.

For the Virasoro shadow tower the invariants are the Stokes pole
\[
 c_S=-\frac{218}{45},
\]
the Borel radius
\[
 |\omega|^2(c)=\frac{c^2(5c+22)}{4(45c+218)},
\]
and the Riccati equation
\[
 U'+2U^2=\frac{Q_c''}{2Q_c},\qquad U=(\log H)'-\frac2t .
\]
They are compared with the conductor splitting
\((K,K^\kappa,K_c,K_g)\), the \(p=691\) Chenevier determinant, and
the Gritsenko--Cl\'ery eight-form spread.

For the chiral-topological theory the comparison passes through the
\(\mathsf{SC}^{\mathrm{ch,top}}\) heptagon and the finite-jet
Khan--Zeng sector.  The shifted-cotangent BV action and the
CME/PVA-Hamiltonian equivalence are constructed.  The all-loop QME,
the reflected weight identity, the boundary vertex algebra, and the
\(E_3\)-lift remain tied to the endpoint hypotheses: the analytic
spectral-decomposition representation, the Fulton--MacPherson Stokes
datum, the reflected weights, and the \(T\)-lift.

For the Calabi--Yau side the comparison uses the two-stage
factorisation
\[
 \Phi_d=\mathrm{Sp}^{\mathrm{ch}}_{\Sigma_{d-1},C}\circ
 \Phi^{\mathrm{FA}}_d,
\]
the K3 Hall--Drinfeld chiral bialgebra
\(\mathbf{H}_{\Delta_5}\), the Joyce vertex algebra with
Kontsevich--Pandharipande--Soibelman orientation data, MNOP
invariants, \(S^3\)-framing on generic compact Calabi--Yau
threefold Lagrangians, and the
\(\mathfrak{so}(4,20)\) K3 Yangian representation-target trichotomy.

The arithmetic branch is governed by the post-Tate spectral-flow
problem over \(\overline{\mathrm{Spec}\,\mathbb Z}\).  The ordinary
\(M_t\)-graph dual, with Hilbert norm
\[
 \|f\|^2=\int_{\mathbb R}(1+t^2)|f(t)|^2\,dt,
\]
has continuous spectral measure and therefore cannot by itself carry
the Hadamard residues on
\[
 \Gamma_\xi^{\mathrm{crit}}
 =
 \{\gamma\in\mathbb R:\xi(\tfrac12+i\gamma)=0\}.
\]
Any spectral-flow construction which survives this obstruction must
live in a stronger boundary topology: a rigged Hilbert space, a
Paley--Wiener or de Branges space, a reproducing-kernel space, a
nonlocal \(L^2\) completion, or a BV--BFV boundary category on which
the Hadamard trace and the jet maps are continuous.

The basic examples fix the scale.  The Borcherds identity above is
the denominator identity, not arithmetic on a table.  The non-critical
affine Kac--Moody central charge
\[
 \kappa=\frac{\dim(\mathfrak g)(k+h^\vee)}{2h^\vee}
\]
is the Feigin--Frenkel screening-cohomology calculation, not a
second reading of the Sugawara character.  The class \(M\) MC5 Wick
coefficient
\[
 S_5=-\frac{48}{c^2(5c+22)}
\]
is the Virasoro OPE contraction coefficient.
